\begin{frame}{Advanced Octave usage - Packages (1/3)}
\begin{columns}
\begin{column}[t]{0.2\textwidth}
\includegraphics[width=\textwidth]{res/libreoffice/octave_c_cpp_fortran_1}
\end{column}
\begin{column}{0.2\textwidth}
\ttfamily\scriptsize
\colorbox{green!30}{nice-package} \\
├ \colorbox{orange!30}{doc} \\
│ └─ html-files \\
├ \colorbox{orange!30}{inst} \\
│ ├─ \colorbox{orange!30}{private} \\
│ │ \ \ └─ m-files \\
│ └─ m-files \\
├ \colorbox{orange!30}{src} \\
│ ├─ C/C++-files \\
│ ├─ configure \\
│ └─ Makefile \\
├ CITATION \\
├ {\color{red!50!black}COPYING}  \\
├ {\color{red!50!black}DESCRIPTION} \\
├ INDEX    \\
├ Makefile \\
├ NEWS     \\
├ PKG\_ADD \\
└ PKG\_DEL \\
\end{column}
\begin{column}{0.5\textwidth}
\color{green!50!black}\scriptsize
package name \\[2em]
manual for the package \\[8em]
called if present \\
called by '\texttt{make}' if present \\
called by '\texttt{citation ("nice-package")}' \\
{\color{red!50!black}required} license \\
{\color{red!50!black}required} \\[1em]
for maintenance \\
list of code changes/improvements \\
commands run, when added   to   users path \\
commands run, when removed from users path
\end{column}
\end{columns}
\end{frame}


\begin{frame}{Advanced Octave usage - Packages (2/3)}
\begin{columns}
\begin{column}[t]{0.2\textwidth}
\includegraphics[width=\textwidth]{res/libreoffice/octave_c_cpp_fortran_1}
\end{column}
\begin{column}{0.2\textwidth}
\ttfamily\scriptsize
\colorbox{green!30}{nice-package} \\
├ \colorbox{orange!30}{doc} \\
│ └─ html-files \\
├ \colorbox{orange!30}{inst} \\
│ ├─ \colorbox{orange!30}{private} \\
│ │ \ \ └─ m-files \\
│ └─ m-files \\
├ \colorbox{orange!30}{src} \\
│ ├─ C/C++-files \\
│ ├─ configure \\
│ └─ Makefile \\
├ CITATION \\
├ {\color{red!50!black}COPYING}  \\
├ {\color{red!50!black}DESCRIPTION} \\
├ INDEX    \\
├ Makefile \\
├ NEWS     \\
├ PKG\_ADD \\
└ PKG\_DEL \\
\end{column}
\begin{column}{0.5\textwidth}
\ttfamily\tiny
\colorbox{red!30}{DESCRIPTION} \\[0.5em]
Name: nice-package \\
Version: 1.0.0 \\
Date: 2019-12-24 \\
Author: First <first@home.net>, Second <first@home.net> \\
Maintainer: First <first@home.net>, Second <first@home.net> \\
Title: Some very nice package \\
Description: A short description of the package.  If this \\
\ description gets too long for one line it can continue \\
\ on the next by adding a space to the beginning of the \\
\ following lines. \\
License: GPLv3+ \\
{\color{red!50!black}Depends: package (>= 1.0.0)} \\
BuildRequires: libopenblas-dev [Ubuntu] openblas-devel \\[2em]

\colorbox{red!30}{INDEX} \\[0.5em]
\# A comment \\
nice-package >$ $> Some very nice package \\
Category Name 1 \\
 function1 function2 function3 \\
 function4 \\
Category Name 2 \\
 function2 function5 \\
\end{column}
\end{columns}
\end{frame}


\begin{frame}{Advanced Octave usage - Packages (3/3)}
\begin{columns}
\begin{column}[t]{0.2\textwidth}
\includegraphics[width=\textwidth]{res/libreoffice/octave_c_cpp_fortran_1}
\end{column}
\begin{column}{0.75\textwidth}
\small
Why all the effort?
\begin{itemize}
\itemsep1em
\item
Grouping of related code.

\item
Users can easily install your package from an \textbf{arbitrary URL}:

\texttt{\scriptsize\color{DarkBlue}
pkg install https://www.url.org/nice-package-v1.0.0.tar.gz}

\begin{figure}
\includegraphics[width=4em]{res/images/octave-flower}
\end{figure}

\item
Upload to \textbf{Octave Forge} \url{https://octave.sourceforge.io/}.
After some package review it can be installed conveniently via:

\texttt{\footnotesize\color{DarkBlue}
pkg install {\color{red!50!black}-forge} nice-package}

\end{itemize}
\end{column}
\end{columns}
\end{frame}
